%! AP Statistics — Unit 4, Lesson 4.10 — Warm-Up KEY
\documentclass[10pt]{article}
\usepackage{apstats-article}
\usepackage{apstats-key}

\begin{document}

\pageheader{Unit 4, Lesson 4.10}{Warm-Up --- Answer Key}
\begin{tabularx}{\linewidth}{X X X}
Name:\;\ans{Key} & Date:\; & Period:\;
\end{tabularx}

\vspace{0.3in}

{\large\bfseries\color{navy} Part 1 \quad Spiral Review}

\vspace{0.12in}

\begin{enumerate}

    \item
    \renewcommand{\arraystretch}{1.8}
    \begin{tabularx}{\linewidth}{@{} >{\bfseries}p{2cm} X X X X @{}}
    $x$ & 0 & 1 & 2 & 3 \\
    \hline
    $P(X=x)$ & 0.10 & 0.35 & \ans{0.30} & 0.25 \\
    \end{tabularx}

    \vspace{8pt}
    Missing probability $= $ \ans{0.30} \hspace{1cm}
    $\mu_X = $ \ans{$0(0.10)+1(0.35)+2(0.30)+3(0.25)=1.70$}

    \vspace{0.18in}

    \item
    \begin{tabularx}{\linewidth}{@{} X X X @{}}
    (a) $\binom{5}{2} = $ \ans{10} &
    (b) $\binom{6}{0} = $ \ans{1} &
    (c) $\binom{4}{4} = $ \ans{1} \\[10pt]
    \end{tabularx}

    \vspace{0.18in}

    \item
    \begin{enumerate}[label=(\alph*), itemsep=6pt]
        \item \ans{$P(A \cap B) = P(A) \cdot P(B) = 0.6 \times 0.4 = 0.24$}
        \item \ans{$P(\text{neither}) = (1-0.6)(1-0.4) = 0.4 \times 0.6 = 0.24$}
    \end{enumerate}

    \vspace{0.18in}

    \item
    \begin{enumerate}[label=(\alph*), itemsep=6pt]
        \item \ans{Yes --- because you replace the marble each time, the composition of the bag is the same for every draw, so draws do not affect one another.}
        \item \ans{$P(\text{all red}) = (0.4)^3 = 0.064$}
    \end{enumerate}

\end{enumerate}

\vspace{0.2in}

{\large\bfseries\color{navy} Part 2 \quad Looking Ahead}

\vspace{0.12in}

\noindent Today we use the structure from Problem 4 --- repeated independent trials
with the same probability --- to build the \textbf{binomial distribution}.

\vspace{0.14in}

\begin{tcolorbox}[colback=greenbg, colframe=greenacc, boxrule=0.8pt, arc=2mm,
    left=4mm, right=4mm, top=3mm, bottom=3mm]
    \small
    \begin{itemize}[leftmargin=1.3em, itemsep=8pt]
  \item A binomial setting requires four conditions (BINS) to be verified from context.
  \item The binomial formula counts both the arrangements and the probabilities.
  \item Always check BINS before using any binomial formula or calculator command.
    \end{itemize}
\end{tcolorbox}

\vspace{0.18in}

\noindent\textcolor{slate}{\small Student responses will vary.}

\end{document}
